\section{Week 11: Monte Carlo Methods and the Metropolis--Hastings
Algorithm}

\begin{quote}
\textbf{Chapter 5: Parameter estimation with data-driven methods (MCMC)}
· Lecture 1 h · Lab 2 h \textbf{Running data:} the Week 2 forgetting
data with the Week 10 priors, so that we can check the sampler against
the grid posterior we already have.
\end{quote}

\begin{center}\rule{0.5\linewidth}{0.5pt}\end{center}

\subsection{Key terms}

Important terms of this lesson, with their Persian equivalents (repeated
even if they appeared in earlier weeks).

{\def\LTcaptype{none} % do not increment counter
\begin{longtable}[]{@{}
  >{\raggedright\arraybackslash}p{(\linewidth - 4\tabcolsep) * \real{0.3333}}
  >{\raggedright\arraybackslash}p{(\linewidth - 4\tabcolsep) * \real{0.3333}}
  >{\raggedright\arraybackslash}p{(\linewidth - 4\tabcolsep) * \real{0.3333}}@{}}
\toprule\noalign{}
\begin{minipage}[b]{\linewidth}\raggedright
English
\end{minipage} & \begin{minipage}[b]{\linewidth}\raggedright
\RL{\CMfa فارسی}
\end{minipage} & \begin{minipage}[b]{\linewidth}\raggedright
\RL{\CMfa توضیح کوتاه}
\end{minipage} \\
\midrule\noalign{}
\endhead
\bottomrule\noalign{}
\endlastfoot
Monte Carlo & Monte Carlo \RL{\CMfa (مونت‌کارلو)} & \RL{\CMfa تقریب با میانگین
نمونه‌های تصادفی} \\
Metropolis algorithm & \RL{\CMfa الگوریتم} Metropolis & \RL{\CMfa پیشنهاد یک حرکت و
پذیرش یا رد آن} \\
parameter & \RL{\CMfa پارامتر} & \RL{\CMfa کمیت قابل تنظیم مدل که بین افراد یا شرایط
فرق می‌کند} \\
estimate / estimator & \RL{\CMfa برآورد} / \RL{\CMfa برآوردگر} & \RL{\CMfa مقدار به‌دست‌آمده
برای پارامتر} / \RL{\CMfa روش به دست آوردن آن} \\
Markov chain Monte Carlo (MCMC) & \RL{\CMfa زنجیره‌ی مارکوف} Monte Carlo &
\RL{\CMfa نمونه‌گیری از پسین با یک زنجیره‌ی تصادفی} \\
forgetting curve & \RL{\CMfa منحنی فراموشی} & \RL{\CMfa کاهش یادآوری با گذشت
زمان} \\
prior & \RL{\CMfa توزیع پیشین} (prior) & \RL{\CMfa باور درباره‌ی پارامتر پیش از دیدن
داده} \\
sampler & \RL{\CMfa نمونه‌گیر} (sampler) & \RL{\CMfa الگوریتمی که از توزیع نمونه
می‌گیرد} \\
grid approximation & \RL{\CMfa تقریب شبکه‌ای} (grid) & \RL{\CMfa محاسبه‌ی پسین روی
شبکه‌ای از مقادیر} \\
posterior & \RL{\CMfa توزیع پسین} (posterior) & \RL{\CMfa باور درباره‌ی پارامتر بعد
از دیدن داده} \\
chain & \RL{\CMfa زنجیره} & \RL{\CMfa دنباله‌ی نمونه‌های} MCMC \\
proposal & \RL{\CMfa پیشنهاد} (proposal) & \RL{\CMfa مقدار جدیدی که الگوریتم امتحان
می‌کند} \\
acceptance rate & \RL{\CMfa نرخ پذیرش} & \RL{\CMfa نسبت پیشنهادهای پذیرفته‌شده} \\
trace plot & \RL{\CMfa نمودار} trace & \RL{\CMfa مقدار نمونه‌ها در برابر شماره‌ی
تکرار} \\
autocorrelation & \RL{\CMfa خودهمبستگی} & \RL{\CMfa وابستگی نمونه‌های پشت سر هم} \\
burn-in / warm-up & \RL{\CMfa دوره‌ی گرم‌شدن} (burn-in) & \RL{\CMfa نمونه‌های اولیه که
دور ریخته می‌شوند} \\
sensitivity (d′) & \RL{\CMfa حساسیت} & \RL{\CMfa فاصله‌ی توزیع علامت و نویز} \\
likelihood & \RL{\CMfa درست‌نمایی} & \RL{\CMfa احتمال داده به‌عنوان تابعی از
پارامترها} \\
starting point & \RL{\CMfa نقطه‌ی شروع} & \RL{\CMfa سوگیری اولیه به سمت یکی از
پاسخ‌ها} \\
response bias & \RL{\CMfa سوگیری پاسخ} & \RL{\CMfa تمایل به یک پاسخ، مستقل از
محرک} \\
large language model (LLM) & \RL{\CMfa مدل بزرگ زبانی} & \RL{\CMfa شبکه‌ی عصبی‌ای که
کلمه‌ی بعدی متن را پیش‌بینی می‌کند} \\
self-consistency & self-consistency \RL{\CMfa (رأی اکثریت)} & \RL{\CMfa چند بار
نمونه‌گیری جواب و انتخاب رأی اکثریت} \\
simulation & \RL{\CMfa شبیه‌سازی} & \RL{\CMfa تولید داده‌ی مصنوعی با اجرای مدل} \\
MCSE & \RL{\CMfa خطای استاندارد} Monte Carlo & \RL{\CMfa خطای شبیه‌سازی در میانگین
پسین} \\
\end{longtable}
}

\begin{center}\rule{0.5\linewidth}{0.5pt}\end{center}

\subsection{Learning outcomes}

By the end of this week, students will be able to:

\begin{enumerate}
\def\labelenumi{\arabic{enumi}.}
\tightlist
\item
  \textbf{Explain} why posterior quantities are integrals, and how Monte
  Carlo sampling turns integrals into averages, with an error that
  shrinks like \(1/\sqrt{S}\).
\item
  \textbf{Describe} the logic of Markov chain Monte Carlo: a chain whose
  long-run distribution is the posterior, using only the
  \emph{unnormalized} posterior.
\item
  \textbf{Implement} the Metropolis (and Metropolis--Hastings) algorithm
  from scratch for a one- and a two-parameter cognitive model.
\item
  \textbf{Tune} the proposal distribution using the acceptance rate,
  trace plots and autocorrelation, and \textbf{identify} burn-in.
\item
  \textbf{Compute} posterior summaries and derived quantities (e.g., a
  half-life) from samples, and \textbf{validate} a sampler against a
  known answer.
\end{enumerate}

\begin{center}\rule{0.5\linewidth}{0.5pt}\end{center}

\subsection{Motivation \& big picture}

Everything we want from a posterior is an \textbf{integral}: the
posterior mean \(\int \theta\, p(\theta\mid\mathbf{y})\,d\theta\), a
marginal posterior (integrating out other parameters),
\(P(d' > 1\mid\mathbf{y})\), the posterior predictive distribution, the
marginal likelihood. In Weeks 9--10 we did these integrals on grids.
That worked for two parameters. A hierarchical SDT model for 20
participants (Week 13) has more than 40 parameters; a grid with only 10
points per parameter would need \(10^{40}\) evaluations.

\textbf{Monte Carlo} methods replace integration by \textbf{sampling}:
if we could draw many values from the posterior, any integral becomes an
average over the draws, and the accuracy depends on the number of draws,
not on the number of parameters. The problem is that we cannot sample
from a posterior directly; we usually cannot even compute its
normalizing constant.

\textbf{Markov chain Monte Carlo (MCMC)} solves this with a beautifully
simple trick: a random walk through parameter space that moves toward
more probable values more often than toward less probable ones, in
exactly the right proportions. After a while, the places the walker
visits are draws from the posterior. The best-known version, the
\textbf{Metropolis algorithm}, fits in 15 lines of Python, and you will
write it today before using any library.

\paragraph{Where is this used?}

\begin{itemize}
\tightlist
\item
  \textbf{All modern Bayesian software} (PyMC, Stan, JAGS, BUGS) is
  built on MCMC. Every hierarchical model in recent psychology and
  neuroscience papers was fitted with some version of it.
\item
  \textbf{Physics and chemistry} (where Metropolis was invented in 1953,
  to simulate particles), \textbf{epidemiology} (estimating infection
  parameters), \textbf{phylogenetics} (reconstructing evolutionary
  trees) and \textbf{finance} all rely on MCMC.
\item
  \textbf{Cognitive science also uses sampling as a \emph{theory of the
  mind}}: some researchers propose that people approximate Bayesian
  inference by mentally sampling a few possibilities, which explains
  certain biases.
\item
  \textbf{AI}: sampling answers from a language model is Monte Carlo
  (see the AI \& LLMs section).
\end{itemize}

\begin{center}\rule{0.5\linewidth}{0.5pt}\end{center}

\subsection{Lecture notes}

\paragraph{In plain words}

Every question we ask a posterior (``what is the average value?'',
``what is the probability that \(d' > 1\)?'') is an integral, and grids
cannot compute integrals over many parameters. \textbf{Monte Carlo}
replaces the integral with an average: if we had many random values
\emph{drawn from} the posterior, the average of those values would be
the posterior mean, the fraction above 1 would be \(P(d' > 1)\), and so
on. The difficulty is getting those draws. The \textbf{Metropolis
algorithm} gets them with a clever random walk. Stand at some parameter
value, propose a small random step, and compare how plausible
(likelihood × prior) the new spot is. If it is more plausible, move
there. If it is less plausible, move there only \emph{sometimes}, with
probability equal to the ratio of plausibilities; otherwise stay where
you are and record the current value again. Run this long enough and the
places you have visited are a sample from the posterior. Remarkably, you
never need the hard normalizing constant, because only ratios are used.

\textbf{Suggested timing (60 min):} Monte Carlo integration 12 · Markov
chains 8 · the Metropolis algorithm 18 · Metropolis--Hastings 5 ·
tuning, burn-in and autocorrelation 12 · preview 5.

\subsubsection{1. Monte Carlo integration}

If \(\theta^{(1)}, \dots, \theta^{(S)}\) are independent draws from
\(p(\theta\mid\mathbf{y})\), then for any function \(g\):

\[
\mathbb{E}[g(\theta)\mid\mathbf{y}] = \int g(\theta)\, p(\theta\mid\mathbf{y})\, d\theta \;\approx\; \frac{1}{S}\sum_{s=1}^{S} g\big(\theta^{(s)}\big).
\]

{\def\LTcaptype{none} % do not increment counter
\begin{longtable}[]{@{}
  >{\raggedright\arraybackslash}p{(\linewidth - 4\tabcolsep) * \real{0.3333}}
  >{\raggedright\arraybackslash}p{(\linewidth - 4\tabcolsep) * \real{0.3333}}
  >{\raggedright\arraybackslash}p{(\linewidth - 4\tabcolsep) * \real{0.3333}}@{}}
\toprule\noalign{}
\begin{minipage}[b]{\linewidth}\raggedright
Want
\end{minipage} & \begin{minipage}[b]{\linewidth}\raggedright
Choose \(g(\theta)\)
\end{minipage} & \begin{minipage}[b]{\linewidth}\raggedright
Monte Carlo estimate
\end{minipage} \\
\midrule\noalign{}
\endhead
\bottomrule\noalign{}
\endlastfoot
Posterior mean & \(\theta\) & Average of the draws \\
\(P(\theta > 1 \mid \mathbf{y})\) & indicator \(\mathbb{1}[\theta > 1]\)
& Proportion of draws above 1 \\
Posterior of a derived quantity, e.g., half-life \(2^{1/d} - 1\) &
\(g(d)\) & Apply \(g\) to each draw, then summarize \\
Credible interval & --- & Percentiles (or HDI) of the draws \\
Marginal posterior of \(\theta_1\) & --- & Simply ignore the other
columns \\
\end{longtable}
}

The \textbf{Monte Carlo standard error} of the mean is
\(\mathrm{SD}(g)/\sqrt{S}\) for independent draws, independent of the
number of parameters. That is why sampling beats grids in high
dimensions.

\paragraph{Reading the math}

\begin{itemize}
\tightlist
\item
  \(\mathbb{E}[g(\theta)\mid\mathbf{y}] = \int g(\theta)\,p(\theta\mid\mathbf{y})\,d\theta\):
  the posterior average of some quantity \(g(\theta)\). Choosing \(g\)
  chooses the question (see the table).
\item
  \(\approx \frac{1}{S}\sum_{s} g(\theta^{(s)})\): replace the integral
  by the plain average over \(S\) draws
  \(\theta^{(1)}, \dots, \theta^{(S)}\). The superscript \((s)\) is the
  draw number.
\item
  \textbf{Error:} \(\mathrm{SD}(g)/\sqrt{S}\). To halve the error you
  need four times as many draws, whatever the number of parameters.
\item
  \textbf{Example:} 10,000 draws from \(\mathrm{Beta}(8,4)\); 9,000 of
  them are above 0.5, so
  \(P(\theta > 0.5 \mid \mathbf{y}) \approx 0.90\), with Monte Carlo
  error \(\sqrt{0.9 \times 0.1/10000} = 0.003\).
\end{itemize}

\subsubsection{2. Markov chains in one paragraph}

A \textbf{Markov chain} is a sequence
\(\theta^{(1)}, \theta^{(2)}, \dots\) in which each value is generated
from the previous one only, by a fixed transition rule. Under mild
conditions, a chain ``forgets'' its starting point and settles into a
\textbf{stationary distribution}: the long-run proportion of time it
spends in each region. MCMC constructs a transition rule whose
stationary distribution is the posterior. Two consequences: - The early
part of the chain still reflects the starting point. It is discarded as
\textbf{burn-in} (warm-up). - Successive draws are \textbf{correlated},
so \(S\) correlated draws carry less information than \(S\) independent
ones (the effective sample size, Week 12).

\subsubsection{3. The Metropolis algorithm}

Let \(\tilde p(\theta) = p(\mathbf{y}\mid\theta)\,p(\theta)\) be the
\textbf{unnormalized} posterior (likelihood × prior). Choose a symmetric
\textbf{proposal distribution}, usually a normal random walk
\(\theta' \sim \mathcal{N}(\theta, \sigma_{\text{prop}})\).

\begin{enumerate}
\def\labelenumi{\arabic{enumi}.}
\tightlist
\item
  Start at some \(\theta^{(0)}\).
\item
  For \(s = 1, 2, \dots, S\):

  \begin{itemize}
  \tightlist
  \item
    Propose
    \(\theta' \sim \mathcal{N}(\theta^{(s-1)}, \sigma_{\text{prop}})\).
  \item
    Compute the acceptance probability \[
    \alpha_{\text{acc}} = \min\left(1,\ \frac{\tilde p(\theta')}{\tilde p(\theta^{(s-1)})}\right).
    \]
  \item
    Draw \(u \sim \mathrm{Uniform}(0, 1)\). If
    \(u < \alpha_{\text{acc}}\), \textbf{accept}:
    \(\theta^{(s)} = \theta'\). Otherwise \textbf{reject}:
    \(\theta^{(s)} = \theta^{(s-1)}\) (the current value is
    \emph{repeated} in the chain).
  \end{itemize}
\end{enumerate}

\textbf{Why it works intuitively:} uphill moves are always accepted;
downhill moves are accepted with probability equal to the density ratio.
So the chain spends time in each region in proportion to its posterior
probability.

\textbf{Why only the unnormalized posterior is needed:} the evidence
\(p(\mathbf{y})\) appears in both numerator and denominator of the ratio
and cancels. This is the key to MCMC's usefulness.

\textbf{Practical details:} - Work in \textbf{logs}: accept if
\(\ln u < \ln\tilde p(\theta') - \ln\tilde p(\theta)\). -
\textbf{Bounded parameters:} if a proposal falls outside the support
(e.g., \(A > 1\)), its prior is 0, so \(\ln\tilde p = -\infty\) and the
proposal is automatically rejected. (Alternatively, sample on a
transformed, unbounded scale, as in Week 3; then the Jacobian must be
added to the log posterior.) - \textbf{Multiple parameters:} propose all
parameters at once with a vector of proposal SDs, or update one
parameter at a time (which leads to Gibbs sampling, Week 12).

\begin{quote}
\textbf{Advanced: detailed balance.} A transition rule
\(T(\theta \to \theta')\) leaves \(\pi\) stationary if
\(\pi(\theta)\,T(\theta \to \theta') = \pi(\theta')\,T(\theta' \to \theta)\)
for all pairs. Check it for Metropolis with
\(\tilde p(\theta') < \tilde p(\theta)\): the left side is
\(\pi(\theta)\, q(\theta'\mid\theta)\, \tilde p(\theta')/\tilde p(\theta)\)
and the right side is \(\pi(\theta')\, q(\theta\mid\theta') \cdot 1\).
With symmetric \(q\) and \(\pi \propto \tilde p\), the two are equal.
\end{quote}

\paragraph{Reading the math}

\begin{itemize}
\tightlist
\item
  \(\tilde p(\theta) = p(\mathbf{y}\mid\theta)\,p(\theta)\): likelihood
  × prior, \emph{without} dividing by the evidence. The tilde
  (\textasciitilde) marks ``unnormalized''.
\item
  \(\alpha_{\text{acc}} = \min(1, \tilde p(\theta')/\tilde p(\theta))\):
  if the proposal is more plausible, the ratio is \textgreater{} 1 and
  \(\min\) caps it at 1 (always accept). If it is half as plausible,
  accept with probability 0.5.
\item
  ``Draw \(u \sim \mathrm{Uniform}(0,1)\); accept if \(u < \alpha\)'' is
  simply a way to make an event happen with probability \(\alpha\) (a
  biased coin).
\item
  \textbf{Worked example, step 3:} the proposal 0.95 is only 0.052 times
  as plausible as 0.80; \(u = 0.30 > 0.052\), so it is rejected and 0.80
  is written into the chain a second time.
\end{itemize}

\subsubsection{4. Metropolis--Hastings: asymmetric proposals}

If the proposal is not symmetric,
\(q(\theta'\mid\theta) \ne q(\theta\mid\theta')\) (e.g., a log-normal
proposal for a positive rate), the acceptance probability needs a
correction:

\[
\alpha_{\text{acc}} = \min\left(1,\ \frac{\tilde p(\theta')\, q(\theta \mid \theta')}{\tilde p(\theta)\, q(\theta' \mid \theta)}\right).
\]

Metropolis is the special case where the \(q\) terms cancel. Gibbs
sampling (Week 12) is another special case, in which every proposal is
accepted.

\subsubsection{5. Tuning and diagnosing a sampler}

The proposal SD \(\sigma_{\text{prop}}\) controls efficiency:

{\def\LTcaptype{none} % do not increment counter
\begin{longtable}[]{@{}
  >{\raggedright\arraybackslash}p{(\linewidth - 6\tabcolsep) * \real{0.2500}}
  >{\raggedright\arraybackslash}p{(\linewidth - 6\tabcolsep) * \real{0.2500}}
  >{\raggedright\arraybackslash}p{(\linewidth - 6\tabcolsep) * \real{0.2500}}
  >{\raggedright\arraybackslash}p{(\linewidth - 6\tabcolsep) * \real{0.2500}}@{}}
\toprule\noalign{}
\begin{minipage}[b]{\linewidth}\raggedright
\(\sigma_{\text{prop}}\)
\end{minipage} & \begin{minipage}[b]{\linewidth}\raggedright
Acceptance rate
\end{minipage} & \begin{minipage}[b]{\linewidth}\raggedright
Behavior
\end{minipage} & \begin{minipage}[b]{\linewidth}\raggedright
Trace plot looks like
\end{minipage} \\
\midrule\noalign{}
\endhead
\bottomrule\noalign{}
\endlastfoot
Too small & Very high (\textgreater{} 0.9) & Tiny steps; explores
slowly; strong autocorrelation & A slowly wandering line \\
Too large & Very low (\textless{} 0.1) & Most proposals land in
low-probability regions and are rejected; the chain sticks & Flat steps
with long plateaus \\
About right & \ensuremath{\approx} 0.2--0.5 & Good mixing & A ``fuzzy caterpillar'' \\
\end{longtable}
}

Classic theoretical results suggest an optimal acceptance rate of about
0.44 for one parameter, falling toward about 0.23 for many parameters.
Treat these as targets, not laws. Tune \(\sigma_{\text{prop}}\) during a
preliminary run (not during the run you keep). Good proposal SDs are of
the order of the posterior SDs, and correlated posteriors need
correlated (or reparameterized) proposals.

\textbf{Burn-in.} Discard the initial part of the chain before it
reaches the high-probability region; the trace plot shows where.
\textbf{Thinning} (keeping every \(m\)-th draw) reduces storage but
throws away information. It is not needed for correct inference.

\textbf{Autocorrelation.} The lag-\(h\) autocorrelation is the
correlation between \(\theta^{(s)}\) and \(\theta^{(s+h)}\). High
autocorrelation means slow exploration and fewer effectively independent
draws.

\textbf{Always validate} a new sampler on a problem with a known answer
(here, the conjugate beta-binomial and the Week 10 grid).

\paragraph{For an engineer}

\begin{itemize}
\tightlist
\item
  \textbf{Monte Carlo integration} is what you do when you estimate a
  BER or an outage probability by simulation: average an indicator over
  many random realizations. The \(1/\sqrt{S}\) error law is the same,
  and it is why low error rates need very long simulations.
\item
  An MCMC chain is a \textbf{Markov process} whose stationary
  distribution is designed to be the posterior; burn-in is the
  transient, the stationary part is the steady state.
\item
  \textbf{Autocorrelation} of the chain is exactly the autocorrelation
  of a random process. Correlated samples carry less information, just
  as correlated noise samples average out more slowly; Week 12 turns
  this into the ``effective sample size''.
\item
  \textbf{Proposal tuning} is a step-size problem like choosing the step
  size of an adaptive filter: too small converges slowly, too large
  becomes unstable (here: most steps are rejected).
\item
  The difference from ordinary simulation: we cannot generate
  independent samples from the target directly, so we build a process
  that converges to it.
\end{itemize}

\paragraph{For the instructor}

\begin{itemize}
\tightlist
\item
  \textbf{Walk the chain on the board.} Draw the curve
  \(\theta^7(1-\theta)^3\) and run the four steps of the worked example
  with a student rolling the ``\(u\)'' values (a random number app
  works). The rejected step that \emph{repeats} 0.80 is the key moment;
  ask why it must be recorded.
\item
  \textbf{An analogy that works:} a tourist exploring a city at night
  who prefers lit streets. Brighter street: always go. Darker street: go
  only sometimes. After a long walk, the time spent on each street is
  proportional to its brightness.
\item
  \textbf{Show the three trace plots} (too small, good, too large) from
  lab Part D side by side. Students learn to read trace plots in five
  minutes this way.
\item
  \textbf{Typical questions.}

  \begin{itemize}
  \tightlist
  \item
    \emph{``Why not just use a fine grid?''} With 40 parameters, a grid
    is impossible; sampling cost does not grow that way.
  \item
    \emph{``How do we know the chain has converged?''} That is Week 12
    (multiple chains, \(\hat R\)).
  \item
    \emph{``Why throw away burn-in?''} Early draws still reflect the
    arbitrary starting point, not the posterior.
  \end{itemize}
\end{itemize}

\begin{center}\rule{0.5\linewidth}{0.5pt}\end{center}

\subsection{Worked example}

\textbf{Model:} recall \(y = 7\) of \(n = 10\) with a flat prior, so
\(\tilde p(\theta) = \theta^7(1-\theta)^3\) (constants dropped). Exact
posterior: \(\mathrm{Beta}(8, 4)\), mean 0.667. Proposal:
\(\theta' = \theta + \mathcal{N}(0, 0.1)\).

{\def\LTcaptype{none} % do not increment counter
\begin{landscape}\small
\begin{longtable}[]{@{}
  >{\raggedright\arraybackslash}p{(\linewidth - 16\tabcolsep) * \real{0.1111}}
  >{\raggedright\arraybackslash}p{(\linewidth - 16\tabcolsep) * \real{0.1111}}
  >{\raggedright\arraybackslash}p{(\linewidth - 16\tabcolsep) * \real{0.1111}}
  >{\raggedright\arraybackslash}p{(\linewidth - 16\tabcolsep) * \real{0.1111}}
  >{\raggedright\arraybackslash}p{(\linewidth - 16\tabcolsep) * \real{0.1111}}
  >{\raggedright\arraybackslash}p{(\linewidth - 16\tabcolsep) * \real{0.1111}}
  >{\raggedright\arraybackslash}p{(\linewidth - 16\tabcolsep) * \real{0.1111}}
  >{\raggedright\arraybackslash}p{(\linewidth - 16\tabcolsep) * \real{0.1111}}
  >{\raggedright\arraybackslash}p{(\linewidth - 16\tabcolsep) * \real{0.1111}}@{}}
\toprule\noalign{}
\begin{minipage}[b]{\linewidth}\raggedright
Step
\end{minipage} & \begin{minipage}[b]{\linewidth}\raggedright
Current \(\theta\)
\end{minipage} & \begin{minipage}[b]{\linewidth}\raggedright
Proposal \(\theta'\)
\end{minipage} & \begin{minipage}[b]{\linewidth}\raggedright
\(\tilde p(\theta)\)
\end{minipage} & \begin{minipage}[b]{\linewidth}\raggedright
\(\tilde p(\theta')\)
\end{minipage} & \begin{minipage}[b]{\linewidth}\raggedright
Ratio
\end{minipage} & \begin{minipage}[b]{\linewidth}\raggedright
\(u\)
\end{minipage} & \begin{minipage}[b]{\linewidth}\raggedright
Decision
\end{minipage} & \begin{minipage}[b]{\linewidth}\raggedright
New \(\theta\)
\end{minipage} \\
\midrule\noalign{}
\endhead
\bottomrule\noalign{}
\endlastfoot
1 & 0.50 & 0.58 & 0.000977 & 0.001636 & 1.675 & 0.42 & accept (uphill) &
0.58 \\
2 & 0.58 & 0.80 & 0.001636 & 0.001678 & 1.026 & 0.71 & accept (uphill) &
0.80 \\
3 & 0.80 & 0.95 & 0.001678 & 0.000087 & 0.052 & 0.30 & \textbf{reject}:
\(u > 0.052\) & 0.80 \\
4 & 0.80 & 0.66 & 0.001678 & 0.002144 & 1.278 & 0.88 & accept (uphill) &
0.66 \\
\end{longtable}
\end{landscape}
}

\textbf{Chain so far:} 0.50, 0.58, 0.80, \textbf{0.80}, 0.66. The
rejected step \emph{repeats} 0.80 in the chain. That is essential:
repeated values are how the chain records that 0.80 is more probable
than 0.95. Dropping them would bias the sample.

\textbf{Interpretation.} In four steps the chain moved from the start
into the region of high posterior density (0.6--0.8) and refused a jump
into the tail. Over thousands of steps, the proportion of time spent
near each value converges to \(\mathrm{Beta}(8, 4)\).

\begin{center}\rule{0.5\linewidth}{0.5pt}\end{center}

\subsection{Lab (2 hours)}

\subsubsection{Goals}

\begin{enumerate}
\def\labelenumi{\arabic{enumi}.}
\tightlist
\item
  Monte Carlo integration with independent draws: see the \(1/\sqrt{S}\)
  error.
\item
  Write a Metropolis sampler and validate it on the beta-binomial model.
\item
  Sample the posterior of the power forgetting model, compare with the
  Week 10 grid, and compute a derived quantity.
\item
  Tune the proposal: compare too-small, good and too-large proposal SDs
  with trace plots, acceptance rates and autocorrelation.
\end{enumerate}

\textbf{Timing:} Part A 15 min · Part B 25 min · Part C 35 min · Part D
30 min · Part E 15 min.

\subsubsection{Setup}

\begin{quote}
\textbf{Code · \RL{\CMfa کد:}} \nolinkurl{code/week11.py}, cell 1 · \RL{\CMfa سلول}
1
\end{quote}

\subsubsection{Part A: Monte Carlo integration (15 min)}

\begin{quote}
\textbf{Code · \RL{\CMfa کد:}} \nolinkurl{code/week11.py}, cell 2 · \RL{\CMfa سلول}
2
\end{quote}

\textbf{Expected output:}

\Needspace{8\baselineskip}
\begin{Verbatim}
Exact: mean = 0.5926, P(theta > 0.65) = 0.1967
S =    100: mean = 0.5865 (MC s.e. 0.0065), P(theta > 0.65) = 0.1800
S =   1000: mean = 0.5935 (MC s.e. 0.0021), P(theta > 0.65) = 0.2160
S =  10000: mean = 0.5937 (MC s.e. 0.0007), P(theta > 0.65) = 0.2022
S = 100000: mean = 0.5927 (MC s.e. 0.0002), P(theta > 0.65) = 0.1971
\end{Verbatim}

Ten times more draws gives about \(\sqrt{10} \approx 3.2\) times less
Monte Carlo error.

\subsubsection{Part B: a Metropolis sampler, validated (25 min)}

\begin{quote}
\textbf{Code · \RL{\CMfa کد:}} \nolinkurl{code/week11.py}, cell 3 · \RL{\CMfa سلول}
3
\end{quote}

\textbf{Expected output:}

\Needspace{6\baselineskip}
\begin{Verbatim}
Acceptance rate: 0.59
Metropolis: mean = 0.5930, SD = 0.0660, 95% interval = [0.458 0.718]
Exact:      mean = 0.5926, SD = 0.0663, 95% interval = [0.46  0.719]
\end{Verbatim}

\subsubsection{Part C: Metropolis for the power forgetting model (35
min)}

Same priors as in Week 10: \(A \sim \mathrm{Beta}(8, 2)\),
\(d \sim \mathrm{Gamma}(2, \text{scale } 0.3)\).

\begin{quote}
\textbf{Code · \RL{\CMfa کد:}} \nolinkurl{code/week11.py}, cell 4 · \RL{\CMfa سلول}
4
\end{quote}

\textbf{Expected output:}

\Needspace{8\baselineskip}
\begin{Verbatim}
Acceptance rate: 0.49
    mean     SD   2.5%  97.5%
A  0.868  0.052  0.762   0.96
d  0.418  0.048  0.325   0.51
Week 10 grid posterior means for comparison: A = 0.867, d = 0.417
\end{Verbatim}

\textbf{Derived quantities} come for free. Compute the posterior of the
half-life \(t_{1/2} = 2^{1/d} - 1\) (Week 2) and the predicted recall
after 60 days:

\begin{quote}
\textbf{Code · \RL{\CMfa کد:}} \nolinkurl{code/week11.py}, cell 5 · \RL{\CMfa سلول}
5
\end{quote}

\textbf{Expected output:}

\Needspace{6\baselineskip}
\begin{Verbatim}
half-life (days)       mean = 4.493, 95% interval = [2.889 7.443]
P(recall) at 60 days   mean = 0.158, 95% interval = [0.113 0.211]
P(half-life < 3 days | data) = 0.039
\end{Verbatim}

\subsubsection{Part D: tuning the proposal (30 min)}

\begin{quote}
\textbf{Code · \RL{\CMfa کد:}} \nolinkurl{code/week11.py}, cell 6 · \RL{\CMfa سلول}
6
\end{quote}

\textbf{Expected output:}

\Needspace{6\baselineskip}
\begin{Verbatim}
too small    proposal SD = [0.002, 0.002]: acceptance = 0.97, lag-10 autocorrelation of d = 0.98, mean of d = 0.473
about right  proposal SD = [0.04, 0.05]: acceptance = 0.50, lag-10 autocorrelation of d = 0.26, mean of d = 0.416
too large    proposal SD = [0.5, 1.5]: acceptance = 0.00, lag-10 autocorrelation of d = 0.94, mean of d = 0.420
\end{Verbatim}

\textbf{What to notice:} the ``too small'' chain accepts almost
everything but crawls (a slowly wandering trace with autocorrelation
near 1 even at lag 60); the ``too large'' chain rejects almost
everything and sits still for long stretches (flat steps). Both give
unreliable estimates from 10 000 iterations: the ``too small'' chain's
mean of \(d\) (0.473) is far from the true posterior mean (0.417), and
the ``too large'' chain's mean happens to be close only because it
barely moved from its starting value (0.42). The well-tuned chain looks
like a fuzzy caterpillar and its autocorrelation decays quickly.

\subsubsection{Part E: interpret (15 min)}

Answer in your notebook: 1. Do the Metropolis posterior summaries for
\((A, d)\) agree with the Week 10 grid? What does that agreement
establish? 2. How would you report the half-life of this participant's
memory, with uncertainty, in one sentence? 3. Why must rejected
proposals be recorded as repeated values in the chain?

\subsubsection{Exercises (increasing difficulty)}

\textbf{Exercise 1 \ensuremath{\star}.} Using the Part B sampler, estimate
\(P(\theta > 0.65\mid\mathbf{y})\) for 30/50 and compare with the exact
value from Part A. How many iterations do you need to get within \ensuremath{\pm}0.01?

\textbf{Exercise 2 \ensuremath{\star}.} Run the Part C sampler from four very different
starting points, e.g., \((0.3, 0.05)\), \((0.99, 2.5)\), \((0.5, 1.0)\)
and \((0.9, 0.1)\). Plot all traces together. After how many iterations
do they become indistinguishable? (This is the idea behind R-hat, Week
12.)

\textbf{Exercise 3 \ensuremath{\star}\ensuremath{\star} (exponential model).} Write
\nolinkurl{log_post_exp} with the Week 10 priors, sample it, and compare
the posterior with the grid. Then compute the posterior predictive
distribution from your samples and reproduce the Week 10 PPC plot.

\textbf{Exercise 4 \ensuremath{\star}\ensuremath{\star} (Metropolis--Hastings with an asymmetric
proposal).} Sample \(d\) with a \textbf{log-normal} random-walk
proposal, \(\ln d' \sim \mathcal{N}(\ln d, 0.1)\), which never proposes
negative values. Derive the Hastings correction
\(q(d\mid d')/q(d'\mid d) = d'/d\) and include it. Show that omitting
the correction biases the posterior.

\textbf{Exercise 5 \ensuremath{\star}\ensuremath{\star}\ensuremath{\star} (Advanced: adaptive tuning).} Write a warm-up
phase that adjusts the proposal SDs every 500 iterations to push the
acceptance rate toward 0.3 (e.g., multiply the SD by 1.2 if the rate is
too high, by 0.8 if too low), then freezes them for the main run. Why
must the adaptation stop before the samples you keep?

\subsubsection{Lab check-off}

Show: (1) the validation of your sampler against the exact Beta
posterior (Part B), (2) the trace plots and summary for the power model
with the comparison to the grid (Part C), and (3) the tuning figure with
acceptance rates (Part D).

\begin{center}\rule{0.5\linewidth}{0.5pt}\end{center}

\subsection{Connection to AI \& LLMs}

\textbf{Sampling answers from an LLM is Monte Carlo.} Because an LLM
samples its output (Week 1 box), asking the same question several times
gives different answers. A popular trick called
\textbf{self-consistency} samples \(k\) answers and takes the majority
vote {[}Wang et al., 2023, ICLR{]}. Does it help, and how much? This is
a Monte Carlo question: we cannot compute the majority-vote accuracy
easily in closed form, but we can simulate it and report the Monte Carlo
standard error, as in this week's lab.

\begin{quote}
\textbf{Code · \RL{\CMfa کد:}} \nolinkurl{code/week11.py}, cell 7 · \RL{\CMfa سلول}
7
\end{quote}

\textbf{Expected output:}

\Needspace{7\baselineskip}
\begin{Verbatim}
k =  1 samples: majority-vote accuracy = 0.406 (MCSE 0.003)
k =  5 samples: majority-vote accuracy = 0.506 (MCSE 0.004)
k = 15 samples: majority-vote accuracy = 0.671 (MCSE 0.003)
k = 45 samples: majority-vote accuracy = 0.873 (MCSE 0.002)
\end{Verbatim}

\textbf{Try it (no API needed).} - \ensuremath{\star} Change \nolinkurl{p_answer} so that
one \emph{wrong} answer is the most likely (e.g.,
\nolinkurl{[0.30, 0.35, 0.20, 0.15]}). What does majority voting do now?
Why? - \ensuremath{\star}\ensuremath{\star} How many simulated questions (\nolinkurl{n_sim}) do you need
for an MCSE of 0.001? Check your answer with the formula
\(\sqrt{p(1-p)/n}\). - \ensuremath{\star}\ensuremath{\star} Connection to cognition: the ``wisdom of the
crowd within'' idea says that averaging a person's repeated guesses
improves accuracy. Which assumption of this simulation (independent
samples) is unlikely to hold for repeated human guesses, or for an LLM
at low temperature?

\begin{center}\rule{0.5\linewidth}{0.5pt}\end{center}

\subsection{Common mistakes \& misconceptions}

\begin{enumerate}
\def\labelenumi{\arabic{enumi}.}
\tightlist
\item
  \textbf{Dropping rejected steps} instead of repeating the current
  value. This biases the sample toward the tails.
\item
  \textbf{Using the normalized posterior.} Not needed and usually
  impossible; the evidence cancels in the ratio.
\item
  \textbf{Computing the ratio without logs.} Likelihoods underflow.
  Compare log densities.
\item
  \textbf{Forgetting the prior support.} Proposals outside the support
  must get \(\ln\tilde p = -\infty\), otherwise the likelihood may be
  evaluated at impossible values (e.g., negative probabilities produce
  \nolinkurl{nan}).
\item
  \textbf{Judging the sampler only by acceptance rate.} A 30\% rate does
  not guarantee convergence. Look at traces and autocorrelation, and run
  several chains (Week 12).
\item
  \textbf{Keeping the burn-in} or tuning the proposal during the samples
  you keep. Both bias the results.
\item
  \textbf{``More iterations always fix it.''} If the chain is stuck in
  one mode, or the proposal is badly scaled, more iterations of the same
  chain may not help in practical time. Diagnose first.
\item
  \textbf{Forgetting the Hastings correction} with asymmetric proposals.
\end{enumerate}

\begin{center}\rule{0.5\linewidth}{0.5pt}\end{center}

\subsection{Check your understanding}

\textbf{1.} Why is the Monte Carlo error independent of the number of
parameters, while grid approximation is not? \textbf{Answer} The MC
estimate is an average of \(S\) draws, and its standard error is
\(\mathrm{SD}/\sqrt{S}\) regardless of the dimension of \(\theta\). A
grid must cover the whole parameter space, so its size grows
exponentially with the number of parameters.

\textbf{2.} In the Metropolis algorithm, why do we not need the evidence
\(p(\mathbf{y})\)? \textbf{Answer} Acceptance depends only on the
\emph{ratio} of posterior densities at the proposed and current values;
\(p(\mathbf{y})\) appears in both and cancels.

\textbf{3.} Current \(\tilde p(\theta) = 0.004\), proposed
\(\tilde p(\theta') = 0.001\), and \(u = 0.18\). Accept or reject? What
if \(u = 0.31\)? \textbf{Answer} The ratio is 0.25. With
\(u = 0.18 < 0.25\), accept; with \(u = 0.31 > 0.25\), reject and repeat
the current value.

\textbf{4.} A chain has acceptance rate 0.97. What is probably wrong,
and how does its trace look? \textbf{Answer} The proposal SD is too
small: almost every tiny step is accepted, the chain explores very
slowly, and the trace looks like a slowly wandering line with high
autocorrelation. Increase the proposal SD.

\textbf{5.} What is burn-in, and how do you decide its length?
\textbf{Answer} The initial iterations, before the chain has reached the
high-probability region, which still reflect the arbitrary starting
value. Choose the length from trace plots (and from several chains
started at different points), and discard those iterations.

\textbf{6.} When is the Hastings correction needed? Give an example.
\textbf{Answer} When the proposal is asymmetric,
\(q(\theta'\mid\theta) \ne q(\theta\mid\theta')\), e.g., a log-normal
proposal for a positive rate, or a proposal that depends on the current
value's position in a bounded interval.

\textbf{7.} How do you obtain the posterior of the half-life
\(t_{1/2} = 2^{1/d} - 1\) from MCMC samples of \(d\)? \textbf{Answer}
Transform every sample of \(d\) into a half-life and summarize the
transformed samples (mean, percentiles, HDI). No extra derivation is
needed; this is one of the big practical advantages of sampling.

\begin{center}\rule{0.5\linewidth}{0.5pt}\end{center}

\subsection{Assignment}

No new graded assignment this week. \textbf{A4 (Bayesian analysis by
grid)} is due at the start of the Week 12 lecture. \textbf{A5 (MCMC from
scratch)} is released in Week 12 and builds directly on today's
\nolinkurl{metropolis} function, so keep your lab notebook tidy.

\textbf{Optional practice (ungraded):} Exercise 2 (multiple starting
points) is the best preparation for Week 12.

\begin{center}\rule{0.5\linewidth}{0.5pt}\end{center}

\subsection{Further reading}

\begin{itemize}
\tightlist
\item
  \textbf{Kruschke (2014), chapter 7.} Markov chain Monte Carlo: an
  intuitive introduction to the Metropolis algorithm (the ``politician''
  example), its logic and representativeness of samples. \emph{(Main
  reading.)}
\item
  \textbf{Farrell \& Lewandowsky (2018), chapter 7.} Bayesian parameter
  estimation with MCMC, applied to cognitive models.
\item
  \textbf{Martin, Kumar \& Lao (2021)}: the sections on inference
  engines and MCMC {[}verify chapter{]}.
\end{itemize}

\emph{Preview for Week 12:} Gibbs sampling (no rejections at all), and
the question every MCMC user must answer: has my chain converged? We
will compute R-hat and the effective sample size ourselves, then check
them against ArviZ.
